\subsection{Two identity traces retaining the actual collision order}
The preceding common-radius result has a genuine multiple-radius
boundary. Both identity traces in the actual End/Hom model can
retune while the full relative order and its tensor action are
preserved.

\begin{finding}[\fid{F54}, Part 2: a genuine trace-only full-order model]
\label{f:traceonlyactualmodel}
There is a finite localization $S$ of $\mathbb Q[v,v^{-1}]$,
admitting $v=1$ and every $v=2^n$, and invertible regular
operators on free modules $H,K$ of ranks $6,44$, such that
\[
 \End_F(H)=A_S,\qquad \operatorname{Hom}_F(H,K)=M_S.
\]
Here $A_S$ is the actual order-two rank-three nonstandard
algebra of \fid{F19}, and $M_S$ is the actual rank-$64$
tensor lattice of \fid{F21}, with its named right generators.
Postcomposition also retains its ordinary left operators.
For every $n\ge1$,
\begin{equation}\label{eq:twogenuinetraces}
 \Tr F_H(2^n)=\Tr(F_H(2)^n),\qquad
 \Tr F_K(2^n)=\Tr(F_K(2)^n).
\end{equation}
The base spectra have integral weights $0,1,2,3$ under
$v=Q^{1/2}$, but full spectral retuning fails.
Thus the two identity traces do not determine the
parameter-dependent spectra or their collision data.
\end{finding}

\begin{proof}
Use
\[
 \delta=v+v^{-1},\quad D=\delta^2,\quad
 \Delta=\frac{(D-1)(D-4)}{D^2},\qquad
 S_0=\mathbb Q[v,v^{-1},\delta^{-1},(D-1)^{-1},(D-2)^{-1}].
\]
The previously proved actual algebra and tensor module are
\[
 A_{S_0}=S_0\times S_0\times M_2(C),\qquad
 C=S_0[x]/(x(x-\Delta)),
\]
\[
 M_{S_0}=\epsilon_+^{20}\oplus\epsilon_-^4
       \oplus N(\delta)^{16}\oplus N(D-2)^4.
\]
The two branch values of $x$ are $0,\Delta$ in that order.
All stated F21 changes of basis are regular over $S_0$.

\par\smallskip\noindent\textbf{Explicit scalar functions.}\quad
Set
\[
 k=\frac{2500}{189},\qquad a=3-v,
\]
\[
 c=-v^3+2v^2+4v-8+k\Delta,\qquad d=c+k\Delta,
\]
\[
 z=6v^3-6v^2-14v+30-6k\Delta.
\]
Collecting terms gives
\begin{align}
 a+z+2c+2d&=1+v+2v^2+2v^3=:S_H(v),
       \label{eq:traceescapeH}\\
 20a+4z+16c+4d&=20+4v+16v^2+4v^3=:S_K(v).
       \label{eq:traceescapeK}
\end{align}
Furthermore,
\[
 \Delta(2)=\frac{189}{625},\quad k\Delta(2)=4,
\]
\begin{equation}\label{eq:traceescapevalues}
 (a,z,c,d)|_{v=2}=(1,2,4,8),\qquad
 (a,z,c,d)|_{v=1}=(2,16,-3,-3).
\end{equation}

\par\smallskip\noindent\textbf{One finite base for all evaluations.}\quad
From $S_0$, invert $a,z,c,d$ and
\[
 a-z,\quad a-c,\quad a-d,\quad z-c,\quad z-d
\]
to obtain $S$.
The difference $d-c=k\Delta$ is not inverted.
Every displayed factor is nonzero at $v=1$ and $v=2$
by \eqref{eq:traceescapevalues}.
For $v>1$, $D>4$, so
\[
 0<\Delta=(1-D^{-1})(1-4D^{-1})<1.
\]
At $v=4$, exact values are
\[
 (a,z,c,d)=
 \left(-1,\frac{17007502}{83521},
       -\frac{1192004}{83521},-\frac{379504}{83521}\right),
\]
giving $c<d<a<0<z$.

For $v\ge8$,
\[
 d-a<-v^3+2v^2+5v-11+2k.
\]
The right side is $-62095/189$ at $8$ and has derivative
$-3v^2+4v+5<0$ on this range.
Thus $c<d<a<0$ there.
For $v\ge4$,
\[
 z>6v^3-6v^2-14v+30-6k.
\]
This lower bound is $11506/63$ at $4$ and has positive
derivative $18v^2-12v-14$ on the range.
Hence $z>0$.
Every factor inverted in $S$ is consequently nonzero
at every $2^n$.
The factors from $S_0$ are also nonzero at $1$ and
every real $v>1$. This proves the finite-base claim.

\par\smallskip\noindent\textbf{The exact algebra and module.}\quad
Let $J=C=S[x]/(x(x-\Delta))$ and give it the operator
\[
 F_J=c+kx
     =\begin{pmatrix}c&0\\ k&d\end{pmatrix}
\]
in the column basis $(1,x)$.
Its determinant $cd$ is a unit, and
$x=(F_J-c)/k$.
Therefore commuting with $F_J$ is the same as being
$C$-linear, and
\[
 \End_{F_J}(J)=C.
\]
This identity persists after every coefficient base change.
At $v=1$, $F_J$ has a genuine nonzero Jordan off-diagonal
entry $k$ and eigenvalue $-3$.

Write $S_r$ for a line with scalar operator $r$, and set
\[
 H=S_a\oplus S_z\oplus J^{\oplus2}.
\]
The scalar difference $a-z$ and the determinants
$(a-c)(a-d)$ and $(z-c)(z-d)$ are units.
They exclude all cross-sector intertwiners.
Consequently
\begin{equation}\label{eq:traceescapefullcentralizer}
 \End_F(H)=S\times S\times M_2(C)=A_S
\end{equation}
over $S$ and after any coefficient base change.
The center is the actual order-two center:
\[
 D-4=(v-1)^2(v+1)^2/v^2
\]
and the remaining factor in $\Delta$ is a unit at $1$.

Let $L_0=C/(x)$ and $L_\Delta=C/(x-\Delta)$, with
operators $c,d$, and put
\[
 K=S_a^{20}\oplus S_z^4\oplus L_0^{16}\oplus L_\Delta^4.
\]
The cyclic source gives
\[
 \operatorname{Hom}_F(J^2,L_i)
       =\operatorname{Hom}_C(C^2,L_i)
       =\operatorname{Row}(L_i).
\]
Precomposition supplies the right matrix action.
Thus
\[
 \operatorname{Hom}_F(H,K)=
 \epsilon_+^{20}\oplus\epsilon_-^4
 \oplus\operatorname{Row}(L_0)^{16}
 \oplus\operatorname{Row}(L_\Delta)^4.
\]
The regular \fid{F21} branch identification makes this
the actual tensor lattice with the same named $P_i$.
Its rank is $20+4+32+8=64$.
The named-action identification uses the explicit regular
comparison, not only the dimension.

Over the domain $S$, maps between the distinct branch
targets vanish because $k\Delta$ is a nonzero divisor.
Their multiplicity matrix units give
\[
 \End_F(K)=M_{20}(S)\times M_4(S)\times
          M_{16}(S)\times M_4(S)
          \cong\End_{A_S}(M_S)
\]
by postcomposition.
This retains the actual ordinary left operators.
At collision, the two target branch sectors merge,
and the two endomorphism algebras acquire the same
cross-branch maps. Their dimension becomes
\[
 20^2+4^2+20^2=816.
\]
The invariant Hom has identity conjugation Frobenius,
as it does in every such invariant-Hom construction.

\par\smallskip\noindent\textbf{The two genuine traces.}\quad
At the base $2$, the eigenvalues $1,2,4,8$ have
multiplicities $1,1,2,2$ on $H$ and $20,4,16,4$ on $K$.
Equations \eqref{eq:traceescapeH}--\eqref{eq:traceescapevalues}
give \eqref{eq:twogenuinetraces} for every $n$.
All multiplicities are genuine and positive.

The spectra themselves fail to retune already at $n=2$:
$a(4)=-1$, whereas every eigenvalue of $F_H(2)^2$ is
positive.
At $Q=4$ the base weights are $0,1,2,3$.
This does not impose those weights on the entire family.
For example $a(8)=-5$ has modulus $5$, not $8^w$ for
any integer $w$.
The model therefore violates neither the single-radius
result of Part~1 nor the integral-weight spectral
comparison of \fid{F51}.
\end{proof}

\paragraph{What the example settles.}
The whole actual relative order, its named tensor action,
the ordinary left operators, and two identity-trace systems
can coexist without full spectral retuning.
The example uses multiple spectral radii and does not
supply a pure family with the arithmetic parameter law.
It does not impose marked-operator traces, a path-sign
comparison, or the canonical integral forms of VIII.
The next part tests those additional trace requirements.

% Separate paragraph for Part III.
\paragraph{\fid{F54}, Part 2, in Grothendieck's language.}
This is a finite family of modules over the actual
coherent center algebra, equipped with a chosen invertible
linear operator and an exact invariant-Hom action comparison.
The nilpotent center survives the collision even when
the two unmarked identity traces fail to determine it.
The affine module/sheaf framework is EGA~I,
Theorem~1.3.7, p.~86, Corollary~1.3.8, p.~87, and
Theorem~1.4.1, pp.~90--92.
The specified action is in the language of SGA~3, I,
4.7.2--3, with cote~90, scan~17 as the audited action
anchor \cite{EGA,SGA3,Archive}.
These framework references are distinct from the
explicit End/Hom and trace proofs supplied here.

